\providecommand{\Res}[2]{\qty[round-mode=figures,round-precision=3]{\fpeval{#1}}{#2}}
\providecommand{\Val}[1]{\num[round-mode=figures,round-precision=4]{\fpeval{#1}}}
\providecommand{\ResSci}[2]{\qty[round-mode=figures,round-precision=3,exponent-mode=scientific]{\fpeval{#1}}{#2}}
\subsubsection*{Question a) Puissance absorbée}
\[ P_{\text{abs}} = \alpha I A = \num{\v{ab}} \cdot \num{\v{I}} \cdot \num{\v{S}} = \Res{\v{ab}*\v{I}*\v{S}}{W} \]

\subsubsection*{Question b) Température d'équilibre}
\begin{enumerate}[label=\arabic*)]
    \item \textbf{Variables connues et recherchées :} $\alpha I = \Res{\v{ab}*\v{I}}{W/m^2}$, $e = \num{\v{em}}$, $T_e = \qty{\fpeval{\v{Te}+273}}{K}$ ; variable recherchée : $T$
    \item \textbf{Équations pertinentes :} à l'équilibre, la puissance absorbée égale la puissance nette rayonnée :
    \[ \alpha I A = e\sigma A\,(T^4 - T_e^4) \]
    \item \textbf{Isolation de la variable recherchée :}
    \[ T = \left(\frac{\alpha I}{e\sigma} + T_e^4\right)^{1/4} \]
    \item \textbf{Résolution numérique :}
    \[ T = \left(\frac{\Val{\v{ab}*\v{I}}}{\num{\v{em}} \cdot \num{5.67e-8}} + (\num{\fpeval{\v{Te}+273}})^4\right)^{1/4} = \Res{(\v{ab}*\v{I}/(\v{em}*5.67e-8) + (\v{Te}+273)^4)^0.25}{K} \]
    \item \textbf{Réponse :} le capteur atteint \textbf{\Res{(\v{ab}*\v{I}/(\v{em}*5.67e-8) + (\v{Te}+273)^4)^0.25}{K}}, soit \textbf{\qty[round-mode=places,round-precision=0]{\fpeval{(\v{ab}*\v{I}/(\v{em}*5.67e-8) + (\v{Te}+273)^4)^0.25 - 273}}{\degreeCelsius}}.
\end{enumerate}

\subsubsection*{Question c) Surface sélective}
Le rayonnement solaire est surtout visible, alors que le capteur chaud émet dans l'infrarouge. Une forte absorptivité dans le visible capte le maximum d'énergie, et une faible émissivité dans l'infrarouge limite les pertes par rayonnement : le capteur devient plus chaud.
